\section{Implications for human--agent organizations}
\subsection{Comparing capability and organizational change}
The specification makes an otherwise ambiguous comparison operational. A capability-only intervention changes actor profiles while keeping assignment, visibility, authority, and improvement procedures fixed. A workflow intervention changes those arrangements under a fixed procedure. A procedure intervention changes how candidate arrangements are evaluated or retained. A factorial study can compare these interventions and their interactions under the same task population and outcome criterion. The bottleneck example in \cref{sec:foundations} supplies one elementary prediction; the mechanism study shows how the retention and acquisition of evidence alter the value of organizational search.

The distinction matters when AI changes both production volume and the work of supervision. Automation can redistribute monitoring demands \citep{bainbridge}; human and machine learning can become organizationally interdependent \citep{sturm}. In the experiment, the same amount of acquired evidence has different value depending on whether the organization reuses it, discards it, or retains it after the relevant task relationship changes. The conditional implication is to evaluate evidence retention together with the workflow that generates evidence.

The supplementary audit model in Appendix~\ref{app:audit} examines the earlier decision to acquire information at all. Its finite-horizon value-of-information controller can acquire little evidence under a strong prior, while the main study asks what happens to evidence after it has been bought. Together the models distinguish acquisition, retention, and use. Neither an audit count nor a revision count alone measures an organization's capacity to learn.

\subsection{Independence, interaction order, and cross-functional breadth}
Independent agent identities do not establish independent evidence. Shared models, training data, tools, or copied explanations can correlate errors, limiting the value of aggregation \citep{clemen}. Social influence can also change the diversity of judgments \citep{lorenz}. An ROI comparison records both pre-existing covariance assumptions and reveal events. It can then vary model diversity, information access, or communication while keeping the decision rule explicit.

Human-first and agent-first interaction differ through the information available when a person commits an assessment. Private commitment can preserve an additional signal, while early advice can supply knowledge the person lacks. The outcome depends on capability, covariance, reliance, and the aggregation rule. The elementary calculation in Appendix~\ref{app:analytical} shows that apparent gains from exposure can equal a reweighting of private evidence. Human--AI studies of reliance and cognitive forcing motivate additional behavioral measurements \citep{bansal,bucinca}; reproducing their effects requires their task and participant conditions, not just matching a curve with selected parameters.

Cross-functional humans can reduce handoffs, interpret errors across domains, and connect agent outputs to downstream requirements. Organization design and coordination requirements identify the dependencies through which these gains could arise \citep{galbraith,cataldo,puranam2021}. The gain must be compared with lost specialist depth, attention limits, and switching costs. A useful experiment varies human breadth while logging transfer count, rework, review delay, and quality at fixed resources. These variables are representable in the specification, while their empirical magnitudes remain task-dependent.

\subsection{Weak-to-strong evaluation and more capable agents}
Weak-to-strong generalization concerns learning under weaker supervision \citep{burns}. In a human--agent organization, a related difficulty appears when reviewers cannot directly assess the outputs or proposed process changes of a more capable agent. Execution capability, evaluation reliability, and authority then become separate design dimensions. The organization can invest in decomposed tests, independent evidence, selective escalation, or stronger evaluation tools; each choice changes the cost and coverage of its improvement procedure. Work on amplified supervision and scalable oversight provides candidate mechanisms for this setting \citep{amplify,oversight}.

Hypothetical artificial superintelligence would widen the gap between execution and human evaluation capability. The same specification would require independently measurable outcomes and explicit authority for revising evaluation procedures.

\subsection{Validity and next empirical test}
The specification is evaluated through a reference checker, a public-data mapping, and a controlled mechanism instance. This establishes a path from named modeling fields to executable decisions and observable omissions. It does not yet establish lower modeling effort, better inter-rater agreement, or better predictions than alternative modeling notations. Those outcomes require users to model the same organizations under competing specifications.

For the mechanism, the catalog, known stratum weights, synthetic error distributions, and cost schedule define the comparison. Program trials evaluate a fixed catalog at a smaller screening budget than operational deployment. Historical program scores average trials conducted at different knowledge states; they are a heuristic selection rule, not a stationary posterior for program quality. Reset discards reusable evidence; cumulative retention assumes past outcomes remain relevant; the fixed window bounds age without detecting a change. These choices explain the limits of the tested learning rules. Concept-drift, switching-bandit, and change-point methods offer principled alternatives for future instances \citep{drift,switchbandit,bocpd}, with acquisition costs represented explicitly \citep{costbandit}.

The acquisition-matched control shows why an improved organizational outcome should not be attributed to the most conspicuous changed component. Periodic program trials also alter when evidence arrives and who can reuse it. The observed gains do not establish that program replacement supplies a stable additional benefit across the tested change timings. A positive result for open-ended procedure discovery would require a broader candidate space and a comparator with equally informative observations.

A direct organizational test would randomize review-routing changes and evaluation coverage across comparable task streams, retain failed proposals, and measure escaped errors and human attention over a declared horizon. A second stage could randomize whether teams may replace the evaluation program. Such a design would separate the value of broader evidence from the value of the procedure used to discover and retain that evidence policy.

\section{Conclusion}
Recursive organization improvement asks how human--agent teams change their work and learn whether those changes helped. The proposed modeling specification connects evidence access, organizational memory, authority, costs, and retention through explicit change contracts. Its mechanism study shows why these connections matter: accumulated evidence largely removes the reset design's penalty for repeated assessment in a stable environment, while obsolete evidence can delay a necessary workflow change. Matching acquisition and reuse then narrows the gain attributable to replacing the evaluator itself. Organizational improvement must therefore be evaluated across the lifecycle of its evidence---what is discovered, what is retained, and when it is reconsidered---rather than inferred from agent capability or the number of revisions performed.
